%%%%%%%%%%%%%%%%%%%%%%%%%%%%%%%%%
%
%       EXERCÍCIO
%
%%%%%%%%%%%%%%%%%%%%%%%%%%%%%%%%%

\definevalorquestao

\begin{exercicioBanco}[\valorquestao]
O quadro a seguir apresenta a quantidade de um tipo de pão vendido em uma semana em uma padaria.

\begin{longtable}{l c}
\hline
Dia da semana & Número de pães vendidos \\
\hline
Domingo       & 250 \\
Segunda-feira & 208 \\
Terça-feira   & 215 \\
Quarta-feira  & 251 \\
Quinta-feira  & 187 \\
Sexta-feira   & 187 \\
Sábado        & 186 \\
\hline
\end{longtable}

O dono da padaria decidiu que, na semana seguinte, a produção diária desse tipo de pão seria igual ao número de pães vendidos no dia da semana em que tal quantidade foi a mais próxima da média das quantidades vendidas na semana.

Qual o dia da semana utilizado como referência para a quantidade de pães a serem produzidos diariamente?
\end{exercicioBanco}

